% 11-transparency.tex: наблюдать за наблюдателями.
\section{Наблюдение за наблюдателями: журналы, свидетели, расщеплённый обзор и публичный отчёт}
\label{sec:transparency}

\subsection{Пробел, остающийся после заякоривания}

Полярис время от времени фиксирует пакет строк аудита в корне Меркла и записывает корень в
таблицу, работающую только на добавление, так что внешний проверяющий способен доказать, что
история под корнем с тех пор не пересматривалась. Один пробел при этом остаётся. Отдельный корень
доказывает лежащие под ним строки; он не доказывает, что оператор позже не отбросил корень и не
опубликовал на его месте другой. Журнал прозрачности закрывает пробел, фиксируя всю
последовательность корней и позволяя всякому доказать, что последовательность только росла.
\sImpl

\subsection{Журнал}

Записи журнала суть корни якорей в порядке добавления; ничего нового не хранится, потому что
лежащая под ними таблица уже работает только на добавление. Над этими записями приложение строит
дерево Меркла в духе Прозрачности Сертификатов (РФС 6962), с ША3-256 и различными префиксами для
листьев и внутренних узлов, так что одно никогда нельзя выдать за другое. Оно публикует
подписанную вершину дерева, то есть обязательство журнала обо всей его истории на некотором
размере; доказательство согласованности того, что дерево размера \textit{м} есть префикс дерева размера
\textit{н}, которое существует лишь в случае, если первые \textit{м} записей ни разу не переставлялись, не
переписывались и не выбрасывались; доказательство включения для любой записи; и сами записи, чтобы
наблюдатель мог воспроизвести журнал и пересчитать всё. Та же механика публикует ещё два журнала:
хеши вычеканенных расписок обмена и хеши заякоренных отметок времени. Личных данных нет ни в чём из
этого. \sImpl

\subsection{Наблюдатель и то, чего он не видит}

Наблюдатель есть самостоятельная программа, которую запускает всякий желающий. Он держит один
кусочек состояния, последнюю принятую вершину; каждый цикл он забирает текущую вершину, сверяет
подпись с ключом, полученным внеполосно, забирает доказательство согласованности от своей последней
вершины до новой и либо записывает новую вершину, либо поднимает тревогу о переписанной записи,
ответвлении, усохшем дереве, вершине, подписанной не тем ключом, или отметке времени, пошедшей
вспять. Доказывает он согласованность: оператор не способен переписать историю так, чтобы этого не
обнаружил всякий наблюдатель, видевший старую вершину. Чего он не способен доказать сам по себе,
так это что двум наблюдателям показали одну и ту же вершину на одном и том же размере. Журнал
способен показать одному наблюдателю частное ответвление истории. \sImpl

% Рисунок открывает подраздел под его заголовком, и заголовок с рисунком остаются на одной
% странице: нужное им вместе место резервируется до заголовка, так что страница, которая не
% вмещает обоих, разрывается перед заголовком, а не между ними.
\Needspace{40\baselineskip}
\subsection{Свидетели, обмен слухами и доказательство двоесказательства}

\figfile{transparency}{Журнал, наблюдатель, свидетели, общий котёл слухов и доказательство
двоесказательства. Журнал публикует подписанную вершину и доказательства. Наблюдатель доказывает,
что журнал только дописывал. Свидетели соподписывают вершины, которые сочли согласованными,
доверяющая сторона требует их порогового числа на точно этой вершине, а свидетели пускают по кругу
те вершины, которые им показали; две подписанные журналом вершины одного размера с разными корнями
суть неотрицаемое доказательство того, что журнал солгал. Вершины можно также публиковать во
внешнюю книгу, работающую только на добавление, которая возвращает расписку.}{fig:transparency}

Свидетель есть наблюдатель, который делает ещё две вещи. Он соподписывает вершину собственным
ключом, и только тогда, когда вершина есть добавляющее продолжение последней им соподписанной, и
он пускает по общему котлу ту подписанную журналом вершину, которую видел. Отсюда следуют два
применения, оба проверяемые автономно. Засвидетельствованная контрольная точка есть вершина,
несущая соподписи по меньшей мере \textit{К} различных доверенных свидетелей, так что расщеплённый обзор
требует двоесказательства от \textit{К} независимых свидетелей, а не от одного лишь журнала.
Доказательство двоесказательства есть две вершины одного журнала, обе верно подписанные ключом
журнала, на одном размере и с разными корнями; журнал подписал обе, а потому отпереться от лжи не
может, и два обменивающихся слухами свидетеля производят это доказательство в тот же миг, как
сверят записи. Свидетель, которому самому показали вторую вершину на размере, им уже
соподписанном, отказывает и пишет доказательство. \sImpl

Внешняя книга довершает картину: журнал публикует каждую вершину в независимую книгу, работающую
только на добавление, и получает расписку, то есть собственную подписанную вершину книги вместе с
доказательством включения, так что журнал не способен воспользоваться вершиной, которую публично не
зафиксировал, а книга не способна позже её выбросить. Книга сама есть журнал только на добавление,
так что здесь переиспользуется та же механика и наблюдают за ней так же. Книга на файловой основе
реализована и проверена тестами; драйверы на основе цепей объявлены и ждут своих интерфейсов.
\sImpl{} для файловой подсистемы; \sFut{} для цепи.

\subsection{Публичный отчёт: сказать вслух, как часто применялась эта власть}

Журналы доказывают, что история не переписывалась. Они не говорят обществу, что история содержит, а
одна её часть важнее прочих: идентификационный орган просит доверить ему власть вытянуть всю
историю проверок одного названного человека по ордеру, и всякий надзор над этой властью, ордер,
проверка роли, модель раскрытия, снаружи невидим. Начиная с \code{в9.382} орган способен
публиковать отчёт о прозрачности, говорящий, как часто он этой властью пользовался, и построенный
из записи о доступе к аудиту, которую такое чтение теперь оставляет (раздел~\ref{sec:privacy}).
Построение отчёта началось с вопроса, что вообще поддаётся записи, и ответом было <<ничто>>,
потому что чтение не оставляло следа о себе. Малые значения не публикуются, а удерживаются, вместо
того чтобы печататься нулём, потому что отчёт с накопительным итогом позволяет читателю вычесть
один период из следующего и восстановить ячейку; поэтому отчёт не несёт бегущего итога, из которого
можно вычитать; а учение при каждой отправке предпринимает ровно такое восстановление и обязано
потерпеть неудачу. Отчёт есть поверхность, на которой орган делает собственный самый навязчивый
поступок исчислимым для посторонних. О доступе, который не был записан, он не говорит ничего, и
пакет для рецензента называет это ограничением собственного журнала любого органа. \sImpl

\subsection{Кто запускает свидетелей}

Сила этой защиты на деле есть в точности независимость тех, кто запускает свидетелей. Построен
протокол: формат соподписи, пороговое правило, доказательство двоесказательства и демон свидетеля,
который может запустить всякий, и всё это отрабатывается в учениях под настоящими подписями при
каждом выпуске, где три свидетеля ловят ответвление двумя способами. Чего не может существовать до
развёртывания, так это набора институционально независимых свидетелей, банка, общественной
организации, другого органа, которые в самом деле всё это запускают. Сегодня свидетели суть три
процесса в учении. Репозиторий говорит, что это не изъян программ, а отсутствие развёртывания, и
настоящий документ это повторяет, потому что слой прозрачности, свидетелей которого запускает та
самая сторона, за которой он наблюдает, слоем прозрачности не является. \sExt

\begin{layercard}{Карточка слоя: прозрачность}
\qa{Что это}{Три журнала в духе РФС 6962 над таблицами только на добавление; подписанные вершины; доказательства согласованности и включения; наблюдатель; свидетели с соподписями, порогом и обменом слухами; доказательство двоесказательства; публикация во внешнюю книгу с расписками; публичный отчёт о собственных чтениях органа по ордеру.}
\qa{Зачем оно существует}{Подписанный объект доказывает, что было сказано; лишь наблюдаемая последовательность доказывает, что сказанное позже не расказали обратно; лишь сверка записей доказывает, что всем сказали одинаково; и лишь опубликованное число позволяет обществу спросить, как часто наблюдаемое случалось.}
\qa{Чему доверяет}{Ключу журнала ради вершины; ключам свидетелей и их независимости, которая есть факт развёртывания; собственной записи органа о своих чтениях.}
\qa{Чему отказывает}{Вершине без доказательства согласованности от предыдущей; вершине, не соподписанной пороговым числом; расписке, которой книга не записала; второй вершине на уже виденном размере; отчёту, из которого малое значение можно восстановить.}
\qa{Что видит}{Корни и хеши, а также числа чтений по ордеру. Никогда человека, никогда полезную нагрузку.}
\qa{Чего не видит}{Сам по себе, расщеплённого обзора: ради этого и нужны свидетели. Доступа, который не был записан.}
\qa{Если откажет}{Переписывание роняет всякого наблюдателя; ответвление есть неотрицаемое доказательство; нечестный свидетель есть один из \textit{К}; выброшенная опубликованная вершина роняет собственную согласованность книги; восстановимое малое значение роняет учение над отчётом.}
\qa{Кем проверяется}{\code{проверка\_журнал\_прозрачности}, \code{проверка\_обмен\_слухами\_прозрачности}, \code{проверка\_публикация\_прозрачности}, \code{проверка\_прозрачность\_отметок\_времени}, \code{проверка\_программа\_прозрачности}, \code{проверка\_поднадзорные\_чтения\_записываются}, оракул канонической равнозначности для подписанной вершины.}
\qa{Сейчас или потом}{Сейчас, как протокол с учением. Независимые свидетели суть развёртывание, потом.}
\qa{Внутри и снаружи}{Внутри: федерация. Снаружи: учреждения, которые обмениваются, подписывают и помечают временем под этим наблюдением.}
\end{layercard}

\nextlayer{Когда полномочия, состояние и история проверяемы посторонними, учреждения могут строить
поверх удостоверения: запрашивать обмен, подписывать документы, привязывать время. Опасность в том,
что ткань обмена станет ровно тем журналом сообщений, который призвание запрещает. Следующий слой
построен как её выворачивание наизнанку.}
